\subsection{新しい進展}
第17章は、数学的基礎が新しい応用分野にも広がっていることを示す。計算神経科学とゲノミクスは、数学モデル、シミュレーション、確率、グラフ、最適化、データ処理を必要とする分野であり、ソフトウェア工学とも深く関係する。

\subsubsection{計算神経科学}
計算神経科学は、神経系の働きを数学モデル、コンピュータシミュレーション、脳の抽象化によって理解しようとする分野である。脳を、情報を符号化し、時間とともに状態が変化する動的システムとして見る。

主な観点には、神経活動が情報をどう符号化するかを扱う神経符号化、単一ニューロンの生物物理、ニューロン同士の相互作用から生じるネットワークの動力学がある。機械学習、人工知能、制御理論、サイバネティクスとも接点がある。

\subsubsection{ゲノミクス}
ゲノミクスは、ゲノムの構造、機能、対応付け、進化、編集を扱う分野である。ゲノムは、DNAにより構成される遺伝情報であり、コンピュータ上で塩基配列を解析することをインシリコ解析と呼ぶ。

ゲノム解析では、DNA配列決定、生命情報学的解析、遺伝的変異の解析、計算モデル化などが使われる。大量の塩基配列データを扱うため、データ管理、セキュリティ、共有、解析効率が重要であり、専用ソフトウェアの開発にはソフトウェア工学が必要である。

\par\smallskip
\begin{center}
\centering
\IfFileExists{assets/generated_figures/ch17/fig-ch17-17.png}{\includegraphics[width=0.96\linewidth]{assets/generated_figures/ch17/fig-ch17-17.png}}{\fbox{\parbox[c][48mm][c]{0.9\linewidth}{\centering 画像生成待ち\\数学的基礎と新しい応用分野の関連図}}}
\par\smallskip
\refstepcounter{figure}\label{fig:ch17-17}図 \thefigure: 数学的基礎と新しい応用分野の関連図
\end{center}
図 \thefigure は、数学的基礎と新しい応用分野の関連図を視覚的に整理したものである。
\par\smallskip
